\chapter{Related Work}
\label{ch:related}

% ─────────────────────────────────────────────────────────────────────────────
\section{Gait-Based Biometric Authentication}
\label{sec:rel_gait_auth}

\subsection{Early Approaches and Feature Engineering}
\label{ssec:rel_gait_early}

Interest in gait as a biometric signal dates to the early 2000s, when
Gafurov et al.~\cite{gafurov2007gait} demonstrated that accelerometer recordings
from waist-mounted devices could distinguish individuals at rates well above
chance using handcrafted cycle-level features. These features proved brittle
when device placement or walking speed varied. The UCI-HAR~\cite{anguita2013human}
and WISDM~\cite{kwapisz2011activity} datasets were originally collected for
activity recognition but became common benchmarks for gait-related tasks due to
their multi-subject, naturalistic recordings.

\subsection{Deep Learning Approaches}
\label{ssec:rel_gait_deep}

Zou et al.~\cite{zou2020deep} showed that a CNN encoder trained directly on raw sensor windows captures person-specific temporal patterns with sufficient discriminability for both identification and pairwise authentication, without requiring explicit step detection. They combine the CNN encoder with an LSTM authentication head trained on pairs of CNN-encoded windows, decoupling representation learning from pairwise comparison. This two-stage architecture is the basis of the system evaluated in this thesis.

\subsection{Privacy in Gait Systems}
\label{ssec:rel_gait_privacy}

Privacy analysis of trained gait models has received considerably less attention
than recognition accuracy. Malekzadeh et al.~\cite{malekzadeh2018mobile} address
gait anonymisation, a different threat model aimed at protecting raw signals
before transmission rather than limiting what a trained model leaks.
Milani~\cite{milani2024gait} analyses identity inference vulnerability in a gait
identification pipeline. The present thesis extends that analysis to an
authentication pipeline and evaluates it across multiple datasets.

% ─────────────────────────────────────────────────────────────────────────────
\section{Membership Inference on Machine Learning Models}
\label{sec:rel_mia_biometric}

\subsection{Foundational Work}
\label{ssec:rel_mia_foundations}

Shokri et al.~\cite{shokri2017membership} formalised membership inference by
training shadow models with known membership labels and using their outputs to
build an attack classifier. Yeom et al.~\cite{yeom2018privacy} grounded the
attack theoretically: the training-test loss gap upper-bounds the attack
advantage, directly connecting memorisation to vulnerability. Salem
et al.~\cite{salem2019ml} showed that a single shadow model and no knowledge
of the target architecture are often sufficient, lowering the bar for practical
attacks. Carlini et al.~\cite{carlini2022membership} unified these ideas into a
likelihood ratio framework (LiRA) that is statistically optimal at any false
positive rate and substantially outperforms earlier methods, particularly on
models with small generalisation gaps; the offline variant is the formulation
adopted in this thesis and is described in detail in
Section~\ref{ssec:bg_lira}.

\subsection{MIA on Biometric Models}
\label{ssec:rel_mia_biometric_detail}

Standard MIA pipelines operate on the softmax output of a classifier, using
the per-class confidence vector as the membership signal. Authentication models
produce no such output: the model scores pairs of windows rather than assigning
class probabilities to individual inputs, so the standard attack surface does
not exist. He and Zhang~\cite{he2021quantifying} studied this problem in the
context of contrastive learning, where models similarly produce embeddings
rather than class probabilities. They adapted existing MIA methods to operate
on the representation space and found that contrastive models are less
vulnerable to membership inference than supervised classifiers, attributing
this to their reduced tendency to overfit individual training examples. This
thesis faces the same structural obstacle and derives an alternative membership
signal from the distribution of pairwise authentication scores, as described
in Chapter~\ref{ch:mia}.

% ─────────────────────────────────────────────────────────────────────────────
\section{Adversarial Attacks on Authentication}
\label{sec:rel_evasion}

\subsection{Adversarial Examples}
\label{ssec:rel_evasion_foundations}

Szegedy et al.~\cite{szegedy2014intriguing} identified that small input
perturbations can reliably change a classifier's output. Goodfellow
et al.~\cite{goodfellow2015explaining} proposed FGSM as a single-step gradient
attack; Madry et al.~\cite{madry2018towards} showed that iterative PGD with
random restart finds stronger examples and serves as a principled robustness
benchmark. The evasion attack in this thesis uses PGD adapted to the pairwise
authentication setting: the goal is an impostor gait window that the model
accepts as the enrolled target.

\subsection{Adversarial Attacks on Biometric Systems}
\label{ssec:rel_evasion_biometric}

Sharif et al.~\cite{sharif2016accessorize} demonstrated physical adversarial
attacks on face recognition via printed eyeglass frames. Gait authentication
presents a different attack surface: the signal is a sensor recording, so
perturbations are applied digitally rather than physically. This thesis
evaluates such digital attacks and assesses whether the required perturbation
budget is consistent with natural within-subject gait variability.

% ─────────────────────────────────────────────────────────────────────────────
\section{Positioning of This Work}
\label{sec:rel_positioning}

This thesis applies membership inference to a pairwise authentication model,
where the standard softmax-based attack surface does not exist and the
membership signal must be derived from the distribution of pair scores. The
evaluation spans multiple datasets under three controlled threat models.
A signal disentanglement experiment isolates the privacy contribution of the
CNN encoder from that of the LSTM authentication head, treating the two
components separately by construction. The evasion evaluation complements
attack success rate with a check against natural within-subject gait
variability, giving a concrete sense of whether successful perturbations
are practically realisable.
